\documentclass[12pt,a4paper]{article}
% Сообщение Ярославу: в работе остались пять моментов. В первой формуле двойного угла для косинуса появился лишний коэффициент; в формуле котангенса неверно записана обратная связь; одно упражнение не доведено до значения квадрата тангенса; в задаче со сдвигом тангенса не подставлено данное числовое значение; последнее упражнение не решено.

\usepackage[a4paper,margin=2cm]{geometry}
\usepackage[T2A]{fontenc}
\usepackage[utf8]{inputenc}
\usepackage[russian]{babel}
\usepackage{amsmath,amssymb,mathtools}
\usepackage{array,booktabs,longtable}
\usepackage{xcolor}
\usepackage[hidelinks]{hyperref}
\usepackage{tikz}

\setlength{\parindent}{0pt}
\setlength{\parskip}{0.6em}
\setlength{\tabcolsep}{3pt}
\renewcommand{\arraystretch}{1.18}

\newcommand{\errorlabel}[1]{\textcolor{red}{\textbf{Ошибка:}} \textbf{#1}}
\newcommand{\keyidea}{\textcolor{blue!60!black}{\textbf{Ключевая идея:}}}
\newcommand{\resultlabel}{\textcolor{teal!60!black}{\textbf{Итог:}}}

\makeatletter
\def\ps@review{%
  \def\@oddhead{}\def\@evenhead{}%
  \def\@oddfoot{\hfil\parbox{0.92\textwidth}{\centering\footnotesize Лёвин Артём Александрович эксклюзивно для Ярослава Гаврилова}\hfil}%
  \let\@evenfoot\@oddfoot
}
\makeatother

\begin{document}

\hypersetup{pageanchor=false}
\begin{titlepage}
\thispagestyle{empty}
\centering
\vspace*{0.23\textheight}
{\Huge\bfseries Ревью домашней работы\par}
\vspace{1.8cm}
{\Large Автор: Лёвин Артём Александрович\par}
\vspace{0.8cm}
{\large 7 октября 2026 года\par}
\vfill
\begin{tikzpicture}[x=1cm,y=1cm]
\draw[line width=0.6pt] (0,0) .. controls (2,0.55) and (4,-0.55) .. (6,0)
                     .. controls (8,0.55) and (10,-0.55) .. (12,0);
\end{tikzpicture}
\end{titlepage}
\clearpage
\hypersetup{pageanchor=true}
\setcounter{page}{1}

\pagestyle{review}
\newpage

\section*{Сводная таблица}
\small
\begin{longtable}{p{1.7cm} p{2.7cm} p{4.3cm} p{3.0cm} p{3.0cm}}
\toprule
\textbf{№ задания} & \textbf{Тема} & \textbf{Суть ошибки} & \textbf{Ваш ответ} & \textbf{Правильный ответ} \\
\midrule
\endfirsthead
\toprule
\textbf{№ задания} & \textbf{Тема} & \textbf{Суть ошибки} & \textbf{Ваш ответ} & \textbf{Правильный ответ} \\
\midrule
\endhead
\hyperref[task:0610f2]{06.10 Ф2} & Формулы двойного угла & В первой из трёх формул перед $\cos^2\alpha$ поставлен лишний коэффициент $2$. & $2\cos^2\alpha-\sin^2\alpha$; две следующие формулы верны. & $\cos^2\alpha-\sin^2\alpha$, $2\cos^2\alpha-1$, $1-2\sin^2\alpha$. \\
\midrule
\hyperref[task:0610f4]{06.10 Ф4} & Котангенс & После верной дроби $\cos\alpha/\sin\alpha$ неверно записана обратная формула. & $\operatorname{ctg}\alpha=\dfrac{\cos\alpha}{\sin\alpha}=-\dfrac{1}{\sin^2\alpha}$. & $\operatorname{ctg}\alpha=\dfrac{\cos\alpha}{\sin\alpha}=\dfrac{1}{\operatorname{tg}\alpha}$, где выражения определены. \\
\midrule
\hyperref[task:0610u3]{06.10 У3} & Основное тригонометрическое тождество & Условие переписано, но вычисление $\operatorname{tg}^2\alpha$ не выполнено. & Ответ не указан. & $\operatorname{tg}^2\alpha=5$. \\
\midrule
\hyperref[task:0710u2]{07.10 У2} & Формулы приведения & Верное преобразование остановлено до подстановки данного значения тангенса. & $-\operatorname{ctg}\alpha$. & $\dfrac{5}{2}$. \\
\midrule
\hyperref[task:0710u3]{07.10 У3} & Выражение через тангенс & Решение и ответ отсутствуют. & Ответ не указан. & $5$. \\
\bottomrule
\end{longtable}
\normalsize

\newpage
\thispagestyle{review}
\vspace*{0.34\textheight}
\begin{center}
{\LARGE\bfseries Разбор ошибок}
\end{center}
\newpage

\phantomsection\label{task:0610f2}
\section*{Задание 06.10 Ф2}

\textbf{Текст задания.} Восстановить три равносильные формулы для $\cos 2\alpha$.

\textbf{Ваше решение.}
\[
\cos 2\alpha=2\cos^2\alpha-\sin^2\alpha,
\qquad
\cos 2\alpha=2\cos^2\alpha-1,
\qquad
\cos 2\alpha=1-2\sin^2\alpha.
\]

\textbf{Ваш ответ.} Первая формула записана как $2\cos^2\alpha-\sin^2\alpha$; две следующие формулы записаны верно.

\errorlabel{лишний коэффициент в первой формуле двойного угла.}

\textbf{Где именно возникает ошибка.} Ошибка появляется сразу в первой из трёх формул. До этого преобразований нет. Первый неверный фрагмент --- коэффициент $2$ перед $\cos^2\alpha$ в записи
\[
2\cos^2\alpha-\sin^2\alpha.
\]
Две последующие записи, $2\cos^2\alpha-1$ и $1-2\sin^2\alpha$, сами по себе верны.

\textbf{Почему этот шаг неверен.} Базовая формула двойного угла для косинуса имеет вид
\[
\cos 2\alpha=\cos^2\alpha-\sin^2\alpha.
\]
Коэффициента $2$ перед $\cos^2\alpha$ здесь нет. Коэффициент $2$ возникает только после использования основного тригонометрического тождества. Например,
\[
\sin^2\alpha=1-\cos^2\alpha,
\]
поэтому
\[
\cos^2\alpha-\sin^2\alpha
=\cos^2\alpha-(1-\cos^2\alpha)
=2\cos^2\alpha-1.
\]
Аналогично, если заменить $\cos^2\alpha$ на $1-\sin^2\alpha$, получается $1-2\sin^2\alpha$. То есть коэффициент $2$ относится к формулам, содержащим только один квадрат тригонометрической функции, но не к исходной разности квадратов.

\keyidea{} Сначала нужно помнить исходную форму $\cos^2\alpha-\sin^2\alpha$, а две остальные получать из неё подстановкой $\sin^2\alpha+\cos^2\alpha=1$.

\textbf{Исправление от первого неверного места.} В первой формуле достаточно убрать лишний коэффициент:
\[
\boxed{\cos 2\alpha=\cos^2\alpha-\sin^2\alpha}.
\]
Две следующие формулы можно оставить без изменений.

\textbf{Полное правильное решение.}
\[
\cos 2\alpha=\cos^2\alpha-\sin^2\alpha.
\]
Из $\sin^2\alpha=1-\cos^2\alpha$ получаем
\[
\cos 2\alpha
=\cos^2\alpha-(1-\cos^2\alpha)
=2\cos^2\alpha-1.
\]
Из $\cos^2\alpha=1-\sin^2\alpha$ получаем
\[
\cos 2\alpha
=(1-\sin^2\alpha)-\sin^2\alpha
=1-2\sin^2\alpha.
\]
Следовательно, три равносильные формулы:
\[
\boxed{\cos 2\alpha=\cos^2\alpha-\sin^2\alpha=2\cos^2\alpha-1=1-2\sin^2\alpha}.
\]

\textbf{Проверка.} При $\alpha=0$ имеем $\cos 0=1$. Все три правильные формулы дают $1$. Ошибочная запись дала бы $2$, поэтому она не может быть тождеством.

\resultlabel{} Первая формула должна быть $\cos^2\alpha-\sin^2\alpha$; две остальные формулы у Вас записаны верно.

\textbf{Задача на закрепление.} Запишите $\cos 2\beta$ тремя равносильными способами: через $\sin\beta$ и $\cos\beta$, только через $\cos^2\beta$ и только через $\sin^2\beta$.

\newpage
\phantomsection\label{task:0610f4}
\section*{Задание 06.10 Ф4}

\textbf{Текст задания.} Восстановить формулу
\[
\operatorname{ctg}\alpha=\frac{\ldots}{\ldots}=\frac{1}{\ldots}.
\]

\textbf{Ваше решение.}
\[
\operatorname{ctg}\alpha=\frac{\cos\alpha}{\sin\alpha}=-\frac{1}{\sin^2\alpha}.
\]

\textbf{Ваш ответ.} $\operatorname{ctg}\alpha=\dfrac{\cos\alpha}{\sin\alpha}=-\dfrac{1}{\sin^2\alpha}$.

\errorlabel{неверная обратная формула для котангенса.}

\textbf{Где именно возникает ошибка.} Последний правильный шаг ---
\[
\operatorname{ctg}\alpha=\frac{\cos\alpha}{\sin\alpha}.
\]
Первый неверный шаг --- переход к
\[
-\frac{1}{\sin^2\alpha}.
\]

\textbf{Почему этот шаг неверен.} Тангенс определяется как
\[
\operatorname{tg}\alpha=\frac{\sin\alpha}{\cos\alpha}.
\]
Если взять обратную величину, получаем
\[
\frac{1}{\operatorname{tg}\alpha}
=\frac{1}{\sin\alpha/\cos\alpha}
=\frac{\cos\alpha}{\sin\alpha}
=\operatorname{ctg}\alpha,
\]
когда все записанные выражения определены. Никакого минуса при обращении дроби не появляется, и квадрат синуса также не возникает.

\keyidea{} Котангенс --- это отношение косинуса к синусу; на общей области определения он является величиной, обратной тангенсу.

\textbf{Исправление от последнего правильного шага.} После
\[
\operatorname{ctg}\alpha=\frac{\cos\alpha}{\sin\alpha}
\]
нужно записать
\[
\boxed{\operatorname{ctg}\alpha=\frac{1}{\operatorname{tg}\alpha}}.
\]

\textbf{Полное правильное решение.}
\[
\operatorname{tg}\alpha=\frac{\sin\alpha}{\cos\alpha}.
\]
Следовательно, на общей области определения
\[
\frac{1}{\operatorname{tg}\alpha}
=\frac{\cos\alpha}{\sin\alpha}.
\]
Поэтому
\[
\boxed{\operatorname{ctg}\alpha=\frac{\cos\alpha}{\sin\alpha}=\frac{1}{\operatorname{tg}\alpha}}.
\]

\textbf{Проверка.} Возьмём $\alpha=\dfrac{\pi}{4}$. Тогда $\sin\alpha=\cos\alpha$, поэтому $\operatorname{ctg}\alpha=1$. Правильная обратная формула тоже даёт $1$, а запись $-1/\sin^2\alpha$ дала бы $-2$.

\resultlabel{} После дроби $\cos\alpha/\sin\alpha$ должно стоять $1/\operatorname{tg}\alpha$ без минуса и без квадрата синуса.

\textbf{Задача на закрепление.} Восстановите пропуски в формуле
\[
\operatorname{ctg}\beta=\frac{\ldots}{\ldots}=\frac{1}{\ldots}.
\]

\newpage
\phantomsection\label{task:0610u3}
\section*{Задание 06.10 У3}

\textbf{Текст задания.} Известно, что
\[
4\sin^2\alpha+10\cos^2\alpha=5.
\]
Найти $\operatorname{tg}^2\alpha$.

\textbf{Ваше решение.} Переписано условие $4\sin^2\alpha+10\cos^2\alpha=5$ и обозначено, что требуется найти $\operatorname{tg}^2\alpha$. Дальнейших преобразований нет.

\textbf{Ваш ответ.} Не указан.

\errorlabel{решение не доведено до вычисления искомой величины.}

\textbf{Где именно возникает ошибка.} Последний правильный шаг --- корректно переписанное условие. Первый отсутствующий шаг --- применение основного тригонометрического тождества, которое позволяет выразить один квадрат через другой.

\textbf{Почему решение неполно.} Само условие ещё не даёт значения $\operatorname{tg}^2\alpha$. Нужно получить отдельно $\sin^2\alpha$ и $\cos^2\alpha$ либо сразу их отношение. Удобнее заменить $\sin^2\alpha$ на $1-\cos^2\alpha$; тогда исходное равенство превращается в обычное линейное уравнение относительно $\cos^2\alpha$.

\keyidea{} В выражении одновременно присутствуют только $\sin^2\alpha$ и $\cos^2\alpha$, поэтому нужно использовать $\sin^2\alpha+\cos^2\alpha=1$, а затем взять отношение $\sin^2\alpha/\cos^2\alpha$.

\textbf{Исправление от последнего правильного шага.} Подставим
\[
\sin^2\alpha=1-\cos^2\alpha.
\]
Тогда
\[
4(1-\cos^2\alpha)+10\cos^2\alpha=5.
\]

\textbf{Полное правильное решение.} Раскрываем скобки:
\[
4-4\cos^2\alpha+10\cos^2\alpha=5.
\]
Собираем подобные слагаемые:
\[
4+6\cos^2\alpha=5.
\]
Отсюда
\[
6\cos^2\alpha=1,
\qquad
\cos^2\alpha=\frac{1}{6}.
\]
По основному тригонометрическому тождеству
\[
\sin^2\alpha=1-\frac{1}{6}=\frac{5}{6}.
\]
Так как $\cos^2\alpha=\dfrac16\ne0$, тангенс определён, и
\[
\operatorname{tg}^2\alpha
=\frac{\sin^2\alpha}{\cos^2\alpha}
=\frac{5/6}{1/6}
=5.
\]

\textbf{Проверка.} Подставим найденные квадраты в исходное равенство:
\[
4\cdot\frac56+10\cdot\frac16
=\frac{20}{6}+\frac{10}{6}
=\frac{30}{6}
=5.
\]
Условие выполнено.

\resultlabel{} $\boxed{\operatorname{tg}^2\alpha=5}$.

\textbf{Задача на закрепление.} Известно, что
\[
3\sin^2\beta+8\cos^2\beta=4.
\]
Найдите $\operatorname{tg}^2\beta$.

\newpage
\phantomsection\label{task:0710u2}
\section*{Задание 07.10 У2}

\textbf{Текст задания.} Известно, что $\operatorname{tg}\alpha=-0{,}4$. Найти
\[
\operatorname{tg}\left(\alpha+\frac{5\pi}{2}\right).
\]

\textbf{Ваше решение.}
\[
\operatorname{tg}\left(\alpha+\frac{5\pi}{2}\right)=-\operatorname{ctg}\alpha.
\]

\textbf{Ваш ответ.} $-\operatorname{ctg}\alpha$.

\errorlabel{верное преобразование не доведено до числового ответа.}

\textbf{Где именно возникает ошибка.} Последний правильный шаг ---
\[
\operatorname{tg}\left(\alpha+\frac{5\pi}{2}\right)=-\operatorname{ctg}\alpha.
\]
Первый отсутствующий шаг --- заменить котангенс величиной, обратной тангенсу, и использовать данное значение $\operatorname{tg}\alpha=-0{,}4$.

\textbf{Почему решение неполно.} В условии дано числовое значение тангенса, поэтому ответ $-\operatorname{ctg}\alpha$ ещё не завершает задачу. Поскольку
\[
\operatorname{ctg}\alpha=\frac{1}{\operatorname{tg}\alpha},
\]
можно сразу получить число. Здесь $\operatorname{tg}\alpha\ne0$, поэтому необходимая обратная величина существует.

\keyidea{} Формула приведения у Вас применена верно; после неё остаётся только выразить котангенс через известный тангенс.

\textbf{Исправление от последнего правильного шага.}
\[
-\operatorname{ctg}\alpha
=-\frac{1}{\operatorname{tg}\alpha}
=-\frac{1}{-0{,}4}.
\]

\textbf{Полное правильное решение.} Так как
\[
\frac{5\pi}{2}=2\pi+\frac{\pi}{2},
\]
то
\[
\operatorname{tg}\left(\alpha+\frac{5\pi}{2}\right)
=\operatorname{tg}\left(\alpha+\frac{\pi}{2}\right)
=-\operatorname{ctg}\alpha.
\]
Далее
\[
-\operatorname{ctg}\alpha
=-\frac{1}{\operatorname{tg}\alpha}
=-\frac{1}{-0{,}4}
=\frac{1}{0{,}4}
=\frac{5}{2}.
\]

\textbf{Проверка.} Поскольку $\operatorname{tg}\alpha=-\dfrac25$, имеем $\operatorname{ctg}\alpha=-\dfrac52$. Тогда $-\operatorname{ctg}\alpha=\dfrac52$, что совпадает с полученным ответом.

\resultlabel{} $\boxed{\operatorname{tg}\left(\alpha+\dfrac{5\pi}{2}\right)=\dfrac52}$.

\textbf{Задача на закрепление.} Известно, что $\operatorname{tg}\beta=\dfrac23$. Найдите
\[
\operatorname{tg}\left(\beta+\frac{3\pi}{2}\right).
\]

\newpage
\phantomsection\label{task:0710u3}
\section*{Задание 07.10 У3}

\textbf{Текст задания.} Известно, что $\operatorname{tg}\alpha=-2{,}5$. Найти значение выражения
\[
\frac{10\cos\alpha+4\sin\alpha+15}{2\sin\alpha+5\cos\alpha+3}.
\]

\textbf{Ваше решение.} Решение отсутствует.

\textbf{Ваш ответ.} Не указан.

\errorlabel{задание не решено.}

\textbf{Где именно возникает ошибка.} Из условия известно значение тангенса, но первый необходимый переход --- связь между синусом и косинусом через тангенс --- не выполнен.

\textbf{Почему решение неполно.} Значение $\operatorname{tg}\alpha=-2{,}5$ означает
\[
\frac{\sin\alpha}{\cos\alpha}=-\frac52.
\]
Тангенс определён, значит $\cos\alpha\ne0$, поэтому можно умножить это равенство на $\cos\alpha$ и получить удобную замену для $\sin\alpha$. После подстановки в исходной дроби слагаемые с $\cos\alpha$ взаимно уничтожаются.

\keyidea{} Не нужно отдельно искать $\sin\alpha$ и $\cos\alpha$. Достаточно использовать отношение $\sin\alpha=\operatorname{tg}\alpha\cos\alpha$.

\textbf{Исправление от первого шага.} Так как
\[
\operatorname{tg}\alpha=-\frac52,
\]
то
\[
\sin\alpha=-\frac52\cos\alpha.
\]
Подставим это выражение в числитель и знаменатель.

\textbf{Полное правильное решение.} Числитель:
\[
10\cos\alpha+4\sin\alpha+15
=10\cos\alpha+4\left(-\frac52\cos\alpha\right)+15.
\]
Так как
\[
4\left(-\frac52\cos\alpha\right)=-10\cos\alpha,
\]
получаем
\[
10\cos\alpha-10\cos\alpha+15=15.
\]
Знаменатель:
\[
2\sin\alpha+5\cos\alpha+3
=2\left(-\frac52\cos\alpha\right)+5\cos\alpha+3.
\]
Следовательно,
\[
-5\cos\alpha+5\cos\alpha+3=3.
\]
Поэтому значение всей дроби равно
\[
\frac{15}{3}=5.
\]

\textbf{Проверка.} После подстановки знаменатель равен $3$, то есть он не обращается в нуль. Следовательно, найденное значение допустимо.

\resultlabel{} $\boxed{5}$.

\textbf{Задача на закрепление.} Известно, что $\operatorname{tg}\beta=-\dfrac32$. Найдите значение выражения
\[
\frac{6\cos\beta+4\sin\beta+14}{4\sin\beta+6\cos\beta+2}.
\]

\end{document}
